Die weichere Feder ist diejenige, die stärker zusammengedrückt wird ($x_\text{w} = \xwO$); die härtere wird um $x_\text{h} = \xhO$ zusammengedrückt. Bei gleicher Kraft gilt
\begin{align}
 F &= D_\text{h}\cdot x_\text{h} = D_\text{w}\cdot x_\text{w}
 \quad\Rightarrow\quad \frac{D_\text{w}}{D_\text{h}} = \frac{x_\text{h}}{x_\text{w}}
\end{align}
Der relative Unterschied der Federkonstante der weicheren zur härteren Feder ist damit
\begin{align}
 \frac{D_\text{w}-D_\text{h}}{D_\text{h}} &= \pDF \\
   &= \frac{\xh}{\xw}-1 \\
   &= \pD \approx \result{\pDP{0}{2}}
\end{align}
Die Federkonstante der weicheren Feder ist also um einen Fünftel kleiner.
